% !TeX root = ../papers/nc-p3.tex
\section{可除代数の特殊化と十五次元成分の量子族との全交叉}
\label{sec:quantum-component-intersection}
\begin{theorem}\label{thm:quantum-component-intersection}
定理\ref{thm:full-fifteen-component}の成分の閉包に属する標準AS窓のうち、
量子多項式代数のindexed窓は添字保持同型を除いて一つだけである。
それは定理\ref{thm:rank-twelve-quantum-limit}の反対階数12の量子窓である。
より一般に、この閉包の標準AS点でFrobenius標準四周期の一つの隅環が
整域なら、その商可除代数の中心上の次数はちょうど4である。
\end{theorem}
\begin{proof}[同じ値群の可除代数の剰余]
まず付値の事実を示す。中心 $F$ 上次数 $d$ の可除代数 $D$ が離散付値を持ち、
値1の中心元 $t$ があるとする。剰余標数を零とする。このとき剰余可除代数も
その中心上の次数が $d$ である。

実際、付値完備化を取ると $D\otimes_F\widehat F\to\widehat D$ は単射である。
左辺は中心的単純環で写像は単位を保つ。像は完備体上有限次元のノルム空間なので完備かつ閉であり、
稠密な $D$ を含むため全射である。従って完備後も次数と中心は指定通りである。
剰余 $\bar D$ の $\bar F$ 上一次独立な元の持上げは、係数の最小付値を除いて還元することにより
$\widehat F$ 上でも独立である。従って剰余基底は有限で、その持上げは逆に全体を張る。
実際、各元の先頭係数を基底で消し、$t$ で割って反復すれば完備性で各係数が収束する。
よって $[\bar D:\bar F]=d^2$ である。

$\alpha\in Z(\bar D)$ は $\bar F$ 上有限分離である。最小多項式を中心の付値環上の
モニックな $f$ に持ち上げる。$\alpha$ の任意の持上げからNewton法を行うと、
各段階はその元の生成する可換部分体の中で行え、導関数の剰余値が非零なので根 $a$ を得る。
この根は非可換な付値環全体でも $\alpha$ の上にただ一つである。
差 $t^m b$ を持つ二根を $t^{m+1}$ を法として比較すると、
$\alpha$ の中心性で一次項は $f'(\alpha)\bar b$ となり矛盾する。
任意の単元による $a$ の共役は同じ剰余の根であり、$a$ は全単元と交換する。
全非零元は $t$ の冪と単元の積なので $a\in\widehat F$ である。従って $Z(\bar D)=\bar F$ を得る。

成分の一般点を分数体側、指定した標準点を特殊側に持つ離散付値環を取る。
例えばその点を吹き上げて正規化した例外素因子で取り、必要なら剰余体を拡張して考える。
基礎体は代数閉体なので、整域性はこの係数拡大で保たれる。
一般点と特殊点で各経路次元は同じ標準値であり、関係は単元な隣接小行列で定まる。
従って各斉次経路加群は自由で、隅の整モデル $A$ の剰余は所望の隅環である。
各元の $t$ 整除回数は有限個の斉次成分で読める。
剰余が整域なのでこの順序は非零元の積で加法的であり、
一般の隅のOre商可除代数に $v(ab^{-1})=v(a)-v(b)$ で延びる。
一般の隅が次数4の可除代数となることは系\ref{cor:quartic-division-corners}による。
値0の分数の分子・分母から共通の $t$ 冪を除けば、剰余は特殊隅の二元の分数となる。
逆に特殊隅の各非零元はこの剰余可除代数内で可逆であり、剰余は特殊隅の商可除代数に一致する。
標準恒等式は係数拡大と降下に可換であるから、剰余体を拡げて得た次数は元の点でも同じである。
これで一般の主張を得た。
\end{proof}

\begin{proof}[量子窓の標準四周期の隅と四次中心]
三角積 $T_{0ij}$ が生成する群の位数を $N$ と書く（無限も許す）。
与えられた量子窓を対角twist $q'_{ij}=q_{ij}d_i/d_j$ で
$\prod_jq'_{ij}=1$ にできる。実際 $d_i^4=(\prod_jq_{ij})^{-1}$ を取り、
その積は四乗根なので一つの $d_i$ の四乗根倍で積を1にする。
この表示では四変数を一度ずつ掛けた体積単項式 $m$ は中心である。
双対量子外積の最高次汎関数について一文字を残りの三文字の後ろへ移したときの係数は
$-\prod_jq'_{ij}$ またはその逆で、全てに共通の符号である。
従って通常の四位置反復が射影基底においてFrobenius標準四周期を与える。

隅は四次Veroneseで、その商可除代数は $\sum_i b_i\equiv0\pmod4$ の
整数指数のLaurent単項式で生成される。実際 $m$ の冪を掛ければ全指数を非負にできる。
この格子は中心な $(1,1,1,1)$ の直線と
$L=\{b:\sum_i b_i=0\}$ の直和である。$L$ の基底 $e_j-e_0$ の交換子は
三角積（または同時にその逆）であり、twistの選択に依らない。
非ねじれ交換子があればその二単項式の量子トーラスはPIにならない。
例えば中心上有限次元の可除代数で一方の共役に対する他方の冪が
無限個の異なる固有値を持つことはできない。
有限の場合、$L$ 上の交代指標の根基の指数は $N^2$ である。
原始 $N$ 乗根に対する三つの指数と $N$ の最大公約数は1であり、
整数の三次交代行列のSmith標準形は非零二因子がともに三係数の最大公約数となる。
従って法 $N$ の像の位数は $N^2$ である（$N=1$ は直接）。
交換子の根基の単項式がLaurent環の中心を張り、その中心分数体上で各剰余類の単項式は基底をなす。
よって隅のPI次数は $N$ である。上の特殊化定理から $N=4$ が必要となる。

さらに成分の四次中心は一般点で五次元なので、標準AS特殊点で少なくとも五次元である。
実際、長さ4,5の経路を自由化すれば矢との交換式の核次元は上半連続である。
量子窓で長さ4の単項式 $x^b$ が中心切断に寄与する必要十分条件は
$b-(1,1,1,1)\in\operatorname{rad}(L)$ である。これは四文字と交換する際の倍率が
文字に依らないことに等しい。その共通倍率は四文字の積 $m$ が中心なので四乗が1であり、
位置別の定数倍を選べば切断が一周する。各多重次数で寄与はちょうど一つである。

三角積を $i^a,i^b,i^c$ と書く。定理\ref{thm:quantum-indexed-all}の全置換で割ると、
位数4の候補と上式による中心次元は以下で尽きる。
\[
\begin{array}{c|ccccc}
 (a,b,c)&(0,0,1)&(0,1,2)&(1,1,1)&(1,1,2)&(1,3,1)\\\hline
 \text{軌道長}&12&24&6&12&2\\
 \text{四次中心次元}&3&2&3&1&5
\end{array}
\]
\path{computations/verify_quantum_component_intersection.py} は原始56組、四添字の全置換と
全次数4単項式を列挙する。残る最後の軌道は定理\ref{thm:rank-twelve-quantum-limit}の
量子窓の三角積であり、定理\ref{thm:walsh-full-component}で実際にこの成分に入る。
従って必要条件がここでは十分である。
\end{proof}
